% Lecture 6: Heapsort & Priority Queues
% (Reminders on matrix multiplication and the heap structure 
% were covered in Lecture 5.)

\subsection{The Heapsort Algorithm}

Heapsort sorts an array \emph{in-place} in $O(n\log n)$ in the worst case.

\begin{enumerate}
  \item \textbf{Build} a max-heap from the input array: $O(n)$.
  \item \textbf{Repeat $n-1$ times:}
    \begin{enumerate}[label=\alph*.]
      \item Exchange the root (maximum) $A[1]$ with the last element $A[heap\_size]$.
      \item Reduce $heap\_size$ by 1 (the exchanged element is in its final position).
      \item Call \textsc{Max-Heapify}$(A, 1)$ to restore the property: $O(\log n)$.
    \end{enumerate}
\end{enumerate}

\begin{center}
\begin{tikzpicture}[font=\small, >=Stealth]
  % Steps illustration
  \node[draw, fill=blue!10, minimum width=1cm, minimum height=0.7cm] (r) at (0,1.5) {16};
  \node[above=2pt of r, font=\tiny, keyred] {max};
  \foreach \v/\i in {14/1, 10/2, 8/3, 7/4, 9/5, 3/6, 2/7, 4/8, 1/9}{
    \node[draw, fill=blue!5, minimum width=0.75cm, minimum height=0.65cm]
          at (\i*0.85-0.5+0.9, 0) {\v};
  }
  \node[draw, fill=green!20, minimum width=0.75cm, minimum height=0.65cm]
        at (0, 0) {16};
  % Arrow
  \draw[->, thick, keyred] (0, 1.1) -- (0, 0.4);
  \node[right=2pt, font=\tiny] at (0, 0.75) {final position};
\end{tikzpicture}
\end{center}

\begin{complexity}{Heapsort}{heapcplx}
\begin{itemize}[noitemsep]
  \item \textsc{Build-Max-Heap}: $O(n)$.
  \item $n-1$ calls to \textsc{Max-Heapify}: $O(\log n)$ each.
  \item \textbf{Total: $O(n \log n)$ in all cases.}
\end{itemize}
\end{complexity}

% ─────────────────────────────────────────────────────────────────────────────
\subsection{Priority Queues}

A \textbf{(max) priority queue} maintains a dynamic set of elements with keys,
and supports:

\begin{center}
\renewcommand{\arraystretch}{1.3}
\begin{tabular}{lll}
\toprule
\textbf{Operation} & \textbf{Description} & \textbf{Time} \\
\midrule
\textsc{Maximum}$(S)$       & Returns the element with the maximum key & $O(1)$ \\
\textsc{Extract-Max}$(S)$   & Removes and returns the maximum      & $O(\log n)$ \\
\textsc{Increase-Key}$(S,x,k)$ & Increases the key of $x$ to $k$   & $O(\log n)$ \\
\textsc{Insert}$(S,x)$      & Inserts $x$ into $S$                 & $O(\log n)$ \\
\bottomrule
\end{tabular}
\end{center}

\textbf{Implementation}: the binary heap implements all these operations efficiently.

\paragraph{\textsc{Extract-Max}.}
\begin{enumerate}[noitemsep]
  \item Save $A[1]$ (the maximum).
  \item Put $A[heap\_size]$ at the root, reduce $heap\_size$.
  \item Call \textsc{Max-Heapify}$(A,1)$.
\end{enumerate}

\paragraph{\textsc{Increase-Key}$(A, i, k)$.}
After updating $A[i] \leftarrow k$, move up the tree as long as the max-heap 
property is violated (the element "floats" upward).

\begin{tcolorbox}[colback=lightblue, colframe=keyblue, fonttitle=\bfseries,
                  title={Application: Priority Queue}]
Used in graph algorithms (Dijkstra, Prim) to efficiently extract the next 
vertex to be processed.
\end{tcolorbox}